%! AP Statistics — Unit 3, Lesson 3.5 — Homework KEY
\documentclass[10pt]{article}
\usepackage{apstats-article}
\usepackage{apstats-key}

\begin{document}

\pageheader{Unit 3, Lesson 3.5}{Homework — KEY}
\namedateperiod

\vspace{0.1in}

\begin{practicebox}
\small
\textbf{1.}

Experimental units: \ans{100 adults}\\[3pt]
Factor: \ans{Sleep aid type}\\[3pt]
Treatments: \ans{New sleep pill; look-alike sugar pill (placebo)}\\[3pt]
Response variable: \ans{Sleep quality rating}\\[3pt]
Design type: \ans{CRD} \quad Blinding: \ans{Double-blind}\\[3pt]
\ans{The placebo controls for the placebo effect — some participants might sleep better simply because they believe they are taking a medication. Without the placebo, we couldn't isolate the drug's true effect.}

\vspace{6pt}
\textbf{2.}

Design type: \ans{CRD (all plants in same greenhouse, no known blocking variable)} \quad Blocking variable: \ans{None identified}\\[3pt]
\ans{If all plants are in the same greenhouse with similar conditions, blocking would not improve the design — there is no nuisance variable to block on. If plants varied by location (e.g., proximity to the window), blocking by location could help.}

\vspace{6pt}
\textbf{3.}

Design type: \ans{Matched pairs design}\\[3pt]
\ans{Current weight likely affects weight-loss outcomes. By pairing on current weight, the researcher removes individual weight differences from the comparison, producing a more precise estimate of the diet effect.}

\vspace{6pt}
\textbf{4.} \ans{Random assignment distributes all potential confounding variables (motivation, health, habits, genetics) roughly equally across the treatment groups by chance. This means any difference in the response variable can be attributed to the treatment rather than to a pre-existing group difference. Without random assignment, a confounding variable could explain the observed difference.}
\end{practicebox}

\end{document}
